Assim como o método básico detalhado no texto [3], a aplicação do Teorema Pi de Buckingham também pode ser descrita como uma receita formal baseada na Seção 2.4.2 do livro [4, 5]. As etapas são:

\begin{enumerate}
    \item[\textbf{a.}] Liste todas as $n$ variáveis derivadas e parâmetros físicos relevantes para o problema [4].
    \item[\textbf{b.}] Identifique as $m$ dimensões primárias (ou fundamentais) envolvidas nessas variáveis (como Massa $M$, Comprimento $L$, Tempo $T$, etc.) [4].
    \item[\textbf{c.}] Calcule o número de produtos ou grupos adimensionais independentes necessários para descrever o sistema, que é dado pela diferença $n - m$ [4].
    \item[\textbf{d.}] Escolha um subconjunto de $m$ variáveis da sua lista original para servirem como \textit{variáveis base} (ou fundamentais). A regra de ouro é que essas $m$ variáveis devem, obrigatoriamente e em conjunto, conter todas as $m$ dimensões primárias, mas não podem de forma alguma formar um grupo adimensional apenas entre si [4].
    \item[\textbf{e.}] Forme os $n - m$ grupos adimensionais ($\Pi_1, \Pi_2, \dots$) permutando cada uma das variáveis restantes uma a uma com as suas $m$ variáveis base. Para isso, multiplique a variável restante pelas variáveis base elevadas a expoentes algébricos desconhecidos ($a, b, c \dots$) [5].
    \item[\textbf{f.}] Para cada grupo formado, force a adimensionalidade aplicando o princípio da homogeneidade dimensional. Substitua as grandezas por suas dimensões primárias e gere um sistema de equações garantindo que o expoente líquido para cada dimensão primária seja zero [6].
    \item[\textbf{g.}] Resolva os sistemas de equações para encontrar os valores dos expoentes ($a, b, c \dots$) e monte os agrupamentos finais [6, 7].
    \item[\textbf{h.}] Expresse o modelo físico final como uma única relação funcional entre os $n - m$ produtos adimensionais independentes (ou seja, $\Pi_1 = \Phi(\Pi_2, \Pi_3, \dots, \Pi_{n-m})$) [8].
\end{enumerate}
